\begin{competences}
  Revoir de manière axiomatique, \cad{} sans rien mémoriser et en sachant retrouver à partir des axiomes, les bases de~:
  \begin{itemize}
  \item Géométrie et son utilisation pour calculer
  \item L'échelle logarithmlique et découvrir son usage pour les calculs
  \item Trigonométrie 
  \item Représentation complexe d'une onde réelle~: phaseurs
  \item Espaces Euclidiens et produits scalaires.
  \end{itemize}
  
\end{competences}

\section{Multiplier et diviser avec Thales}

\question{Calcul de racines carré}

Utilisez le théorème de Pythagore pour contruire une spirale des racines carrées et mesurer les longueurs $\sqrt{2}$, $\sqrt{3}$ et $\sqrt{5}$ sur un graphique. Sachant que
la construction peut être précise au milimètre près ($\pm 0.5mm$), quelle est le
nombre de chiffres significatifs de ces mesures.
Peut-on assurer $10\%$, $5\%$, $1\%$, ou $0.5\%$ d'erreur de calcul graphique ?
Que devient la précision lorsquel'on avance dans la spirale et comment limiter ce problème ?

\question{Précision et Thalès}
Pour une figure de Thales, on peut dire que ~:
\begin{itemize}
\item O/N --- On a forcément un petit triangle dans un grand~;
\item O/N --- Un petit triangle dans un grand peut-être une figure de Thales~;
\item O/N --- Il est nécessaire d'avoir deux côtés perpendiculaires à un troisième côté~;
\item O/N --- Il est nécessaire et suffisant d'avoir un sommet commun et les côtés opposées à ce sommet parallèles entre-eux.
\end{itemize}

Vos réponses à cette question doivent correspondre à un petit robot de \emph{Ulysse 31}.

Utilisez une figure de Thalès pour mesurer $\sqrt{2}$ avec 3 chiffres significatifs et une résolution meilleure que $1\%$.

\question{Règle de trois et Thalès}
Posez votre \og{} produit en croix\fg{} avec votre méthode habituelle pour résoudre le problème suivant~:
\og{} Monsieur Acco pèse maintenant 106 Kg après avoir augmenté sont poids  de $15\%$ pendant cet été. Combien pesait-il au début du mois de Juin ?\fg{}

Calculez à la main le résultat de la division pour avoir 3 chiffres significatifs.

Représentez maintenant ce problème sous la forme d'une figure de Thalès et mesurez le résultat.

Utilisez la bissectrice et faites une construction en escalier pour estimer le poids pour quelques années à venir et les années passées. Que pensez-vous de la rapidité et de la précision de cette méthode ?

\question{L'échelle logarithmique pour multiplier}




%\begin{minipage}{0.4\linewidth}
%  \refWeb{Corrigé \'Echantillonnage 1.1}{https://moodle.insa-toulouse.fr/mod/resource/view.php?id=104716}
%\end{minipage}




%%% Local Variables:
%%% mode: LaTeX
%%% TeX-master: "poly_td"
%%% End:
